% Academic CV. Source layout based on the supplied resume.cls.
\documentclass[margin,line]{resume}
\usepackage{url}
\usepackage{xcolor}
\usepackage{hyperref}
\definecolor{linkblue}{RGB}{23,70,106}
\hypersetup{colorlinks=true,urlcolor=linkblue,pdfauthor={Mukul},pdftitle={Mukul -- Academic CV},pdfsubject={Superconducting quantum hardware and device physics}}
\urlstyle{same}
\input{glyphtounicode}
\pdfgentounicode=1
\setlength{\textheight}{260mm}
\setlength{\paperwidth}{210mm}
\setlength{\paperheight}{297mm}
\pdfpagewidth=\paperwidth
\pdfpageheight=\paperheight
\renewcommand{\labelitemi}{$\bullet$}
\setlength{\leftmargini}{1.4em}
\makeatletter
\g@addto@macro\@listi{\setlength{\itemsep}{2pt}\setlength{\topsep}{3pt}\setlength{\parsep}{0pt}}
\makeatother
\setlength{\parskip}{3pt}
\newcommand{\projectlink}[2]{\href{https://mukul-quantum-researcher.github.io/projects/#1/}{#2}}
\newcommand{\cvheading}[1]{\section{\mysidestyle #1}}

\begin{document}
\name{\Large Mukul}
\begin{resume}

\cvheading{Contact\\Information}
Bengaluru, India \hfill \href{mailto:malikmukul2000@gmail.com}{malikmukul2000@gmail.com}\\
Website: \href{https://mukul-quantum-researcher.github.io/}{mukul-quantum-researcher.github.io}\\
GitHub: \href{https://github.com/Mukul-Quantum-Researcher}{github.com/Mukul-Quantum-Researcher}\\
LinkedIn: \href{https://www.linkedin.com/in/mukul-mukul-47b9b1192/}{mukul-mukul-47b9b1192}

\cvheading{Research\\Profile}
Project Engineer at C-DAC working on superconducting quantum circuits. My experience includes circuit and electromagnetic simulation, cross-resonance gate modelling, readout-filter design and decoherence analysis. I completed an M.Tech at IISc with research in VO$_2$ phase behaviour and coursework in quantum technology. I am interested in PhD research that connects device design, quantum control and experiments on superconducting processors.

\cvheading{Research\\Interests}
Superconducting circuit QED; two-qubit gates and control; readout and calibration; noise and decoherence; parametric amplification. Related interests include semiconductor quantum devices and low-dimensional materials.

\cvheading{Education}
\textbf{Indian Institute of Science, Bengaluru} \hfill \textbf{2023--2025}\\
\textsl{M.Tech in Semiconductor Technology; Minor in Quantum Technology}\\
CGPA: 7.8/10. Thesis supervisor: Prof. Pavan Nukala.\\
Thesis: \projectlink{vo2-thesis}{Stress-Induced Phase Transitions in Vanadium Dioxide (VO$_2$) Nanowires}.\\
Completed 84 credits, including 21 credits in Quantum Technology beyond degree requirements. IISc M.Tech Graduate Fellowship: INR 288,000.

\vspace{1mm}
\textbf{DCR University of Science and Technology, Sonipat} \hfill \textbf{2019--2023}\\
\textsl{B.Tech in Mechanical Engineering}\\
CGPA: 7.75/10.

\cvheading{Research\\Experience}
\textbf{Project Engineer, Quantum Computing Group}\\
\textsl{C-DAC, Bengaluru} \hfill \textbf{July 2025--Present}
\begin{itemize}
\item Contribute to superconducting quantum processor development under India's National Quantum Mission, with work on circuit design, simulation and pulse-level modelling.
\item Compare cross-resonance and parametric gate approaches, examining gate duration, residual ZZ, leakage and coherence trade-offs.
\item Design and simulate qubits, resonators, tunable couplers and Purcell-filter structures using circuit and electromagnetic tools, including Qiskit Metal, Ansys HFSS and Keysight ADS Quantum Pro.
\item Build coupled-transmon CR/ECR simulations, including static-spectrum checks, pulse calibration, leakage analysis and approximate circuit targets for layout. \projectlink{cr-ecr}{Project and code}.
\item Simulate JTWPA gain under different readout-chain terminations in JosephsonCircuits.jl, comparing ideal loads, circulator models and reflectionless filtering. \projectlink{reflectionless-filter}{Project and code}.
\item Model noise and decoherence with open-system methods and study qubit characterisation protocols, including T$_1$/T$_2$, gate-fidelity estimation and randomised benchmarking. Contribute to hardware-aware Max-Cut work and give internal seminars on circuit QED and decoherence.
\end{itemize}

\textbf{Quantum Algorithms Intern, Entangled Ket} \hfill \textbf{Aug--Oct 2024}
\begin{itemize}
\item Implemented variational quantum algorithms with hardware noise models and limited circuit depth; developed educational material on hardware-aware benchmarking.
\end{itemize}

\cvheading{Selected\\Research\\Projects}
\begin{minipage}{\linewidth}
\textbf{Cross-resonance and echoed CR gate simulation} \hfill \projectlink{cr-ecr}{Details} $\vert$ \href{https://github.com/Mukul-Quantum-Researcher/Superconducting-qubits-hardware}{Code}
\begin{itemize}
\item Built exact-cosine static checks and three-level Duffing pulse models for coupled transmons in QuTiP, alongside an earlier scqubits tutorial.
\item Calculated dressed levels and ZZ, checked effective CR rates, compared rotating- and lab-frame dynamics, and calibrated an echoed sequence across a detuning/coupling scan.
\item Selected simulated example: approximately 154 ns ECR duration, 99.685\% coherent average gate fidelity after virtual-Z optimisation and average leakage of about $5.7\times10^{-5}$.
\end{itemize}
\end{minipage}\par\vspace{3pt}

\begin{minipage}{\linewidth}
\textbf{Qubit spectroscopy and Rabi chevrons} \hfill \projectlink{qubit-spectroscopy}{Details} $\vert$ \href{https://github.com/Mukul-Quantum-Researcher/Superconducting-qubits-hardware/blob/main/notebooks/qubit_spectroscopy.ipynb}{Code}
\begin{itemize}
\item Simulated four-level transmon spectroscopy and Rabi dynamics in QuTiP, with relaxation, dephasing, finite-temperature excitation and quasi-static frequency offsets.
\item Fitted frequency--duration chevrons and pulse-duration data; recovered the example's approximately 5.011 GHz transition and 5 MHz Rabi frequency from saved simulated populations.
\end{itemize}
\end{minipage}\par\vspace{3pt}

\begin{minipage}{\linewidth}
\textbf{MTL-coupled Purcell-filter design} \hfill \projectlink{purcell-filter}{Details} $\vert$ \href{https://github.com/Mukul-Quantum-Researcher/purcell-filter-design}{Design files}
\begin{itemize}
\item Implemented a feedline/filter/readout-resonator structure in Qiskit Metal and Ansys HFSS, using a multiconductor transmission-line coupling section to create an intrinsic notch.
\item Iterated shorted-end lengths and examined S/Z-parameter sweeps. The simulations show the notch moving from approximately 11.42 GHz to 8.35 GHz and then 6.83 GHz.
\end{itemize}
\end{minipage}\par\vspace{3pt}

\begin{minipage}{\linewidth}
\textbf{JTWPA gain with reflectionless filtering} \hfill \projectlink{reflectionless-filter}{Details} $\vert$ \href{https://github.com/Mukul-Quantum-Researcher/reflectionless-filter-jtwpa}{Code}
\begin{itemize}
\item Compared one JosephsonCircuits JTWPA under ideal loads, nonideal circulator models and a reflectionless-filter/circulator cascade using harmonic balance and small-signal analysis.
\item Built frequency-dependent termination impedances from surrogate S-parameters and examined gain and external reflection. The study uses model networks rather than the paper's measured Touchstone data.
\end{itemize}
\end{minipage}\par\vspace{3pt}

\begin{minipage}{\linewidth}
\textbf{VO$_2$ phase behaviour and stress-transfer experiments} \hfill \projectlink{vo2-thesis}{Details} $\vert$ \href{https://mukul-quantum-researcher.github.io/assets/documents/Mukul-VO2-Thesis.pdf}{Thesis}
\begin{itemize}
\item Investigated VO$_2$ wires using Raman spectroscopy, stretching on PDMS and AFM nanoindentation during my M.Tech at IISc.
\item Observed width- and substrate-dependent Raman responses under optical excitation. Identified poor strain transfer on PDMS and AFM probe/loading limits; the experiments did not demonstrate a controlled stress-induced transition.
\end{itemize}
\end{minipage}\par\vspace{3pt}

\begin{minipage}{\linewidth}
\textbf{Silicon-on-insulator ring-modulator simulation} \hfill \projectlink{ring-modulator}{Details} $\vert$ \href{https://mukul-quantum-researcher.github.io/assets/documents/Mukul-Ring-Modulator-Report.pdf}{Report}
\begin{itemize}
\item Independently completed a photonics design workflow using Lumerical MODE, FDTD, CHARGE and INTERCONNECT, linking waveguide modes, coupling, PN-junction depletion and ring transmission near 1550 nm.
\item Optimised the passive ring and examined DC-bias-dependent transmission. The passive model gives Q of approximately $5.8\times10^4$ with an FSR near 11 nm.
\end{itemize}
\end{minipage}\par\vspace{3pt}

\cvheading{Other\\Projects}
\begin{itemize}
\item \textbf{Hardware-constrained state preparation:} Five-qubit screening-task solution using X, H, RZ and CX gates on a specified connectivity graph; submitted circuit depth 23. \href{https://github.com/Mukul-Quantum-Researcher/QOSF-Cohort-10}{Notebook}.
\item \textbf{Classiq harmonic-oscillator challenge:} Implemented a quantum differential-equation approach and examined energy estimates and circuit synthesis cost. \href{https://github.com/Mukul-Quantum-Researcher/Q-SITE-Classiq-Open-Challenge}{Submission}.
\item \textbf{Fabrication and characterisation coursework:} Semiconductor pn-junction fabrication using lithography and thin-film processing; graphene characterisation using Raman spectroscopy and AFM.
\end{itemize}

\cvheading{Technical\\Skills}
\textbf{Modelling and programming:} Python, QuTiP, Qiskit, scqubits, NumPy, SciPy, SymPy, Matplotlib, PennyLane; Julia and JosephsonCircuits.jl; MATLAB, Git and \LaTeX.\\
\textbf{Device and EM tools:} Qiskit Metal, Ansys HFSS/Q3D, Keysight ADS Quantum Pro, KLayout, pyAEDT/Elmer and scikit-rf; Ansys Lumerical.\\
\textbf{Experimental background:} Raman spectroscopy, AFM, fabrication coursework, lithography, thin-film deposition and etching.\\
\textbf{Physics:} Circuit QED, coupled-transmon Hamiltonians, pulse dynamics, open quantum systems, semiconductor device physics and integrated photonics.

\cvheading{Selected\\Coursework}
Quantum measurement and sensing; quantum optics; open quantum systems; quantum computation; integrated photonics; quantum transport; semiconductor device physics; micro- and nanofabrication.

\cvheading{Achievements}
\begin{itemize}
\item GATE Mechanical Engineering, 2023: All India Rank 792.
\item Winner, Haiqu Track, Q Site Open Challenge, 2024.
\item 9th place globally, QCoder Programming Contest.
\end{itemize}

\end{resume}
\end{document}
